% Einheit 2 von 5 – Kürzen und Erweitern (bank/brueche-dezimalzahlen/e2.jsonl, zusammenbau v0.1)
%% TODO Zweigzeile Teil 1: was hier gelernt wird (Satz in Schülersprache aus dem Einheitstitel) – Einheit 2
\einheitenkopf[e2]{Einheit 2 von 5 · Kürzen und Erweitern}
\zweigzeile{ab Kl. 5, je nach Buch bis Kl. 6 · P10 oft}
\setcounter{aufgabe}{13}

%% TODO Titel: Ich-kann-Satz fehlt in der Bank (Platzhalter: Kettenname) – Nr. 14
%% TODO Anweisung: ein Satz über der ersten Teilaufgabe fehlt in der Bank – Nr. 14
\begin{aufgabe}{Kürzen und Erweitern}
\begin{teile}
\teil $\frac{2}{5} \to \frac{6}{15}$ – mit welcher Zahl erweitert?
\end{teile}
%% TODO Darstellung für schwach fehlt in der Bank (brueche-dezimalzahlen-e2-k1-s1-v1; Abschnitt „Für schwache Schüler“)
\swz{Erweitere $\frac{3}{5}$ mit $4$.}{}
%% TODO Darstellung für schwach fehlt in der Bank (brueche-dezimalzahlen-e2-k1-s1-v2; Abschnitt „Für schwache Schüler“)
\swz{Erweitere $\frac{1}{6}$ mit $3$.}{}
%% TODO Darstellung für schwach fehlt in der Bank (brueche-dezimalzahlen-e2-k1-s1-v3; Abschnitt „Für schwache Schüler“)
\swz{Erweitere $\frac{4}{9}$ mit $2$.}{}
%% TODO Darstellung für schwach fehlt in der Bank (brueche-dezimalzahlen-e2-k1-s1-v4; Abschnitt „Für schwache Schüler“)
\swz{Erweitere $\frac{5}{7}$ mit $3$.}{}
%% TODO Darstellung für schwach fehlt in der Bank (brueche-dezimalzahlen-e2-k1-s1-v5; Abschnitt „Für schwache Schüler“)
\swz{Erweitere $\frac{2}{3}$ mit $5$.}{}
\swfrage{Was bleibt gleich, was ändert sich?}
%% TODO Darstellung für schwach fehlt in der Bank (brueche-dezimalzahlen-e2-k1-s2-v1; Abschnitt „Für schwache Schüler“)
\swz{Kürze $\frac{14}{21}$ mit $7$.}{}
%% TODO Darstellung für schwach fehlt in der Bank (brueche-dezimalzahlen-e2-k1-s3-v1; Abschnitt „Für schwache Schüler“)
\swz{Kürze $\frac{16}{24}$ vollständig.}{}
%% TODO Darstellung für schwach fehlt in der Bank (brueche-dezimalzahlen-e2-k1-s4-v1; Abschnitt „Für schwache Schüler“)
\swz{Schreibe $\frac{3}{4}$ mit dem Nenner $20$.}{}
%% TODO Darstellung für schwach fehlt in der Bank (brueche-dezimalzahlen-e2-k1-s5-v1; Abschnitt „Für schwache Schüler“)
\swz{Mache gleichnamig: $\frac{1}{3}$ und $\frac{5}{9}$.}{}
%% TODO Darstellung für schwach fehlt in der Bank (brueche-dezimalzahlen-e2-k1-s6-v1; Abschnitt „Für schwache Schüler“)
\swz{Mache gleichnamig: $\frac{3}{4}$ und $\frac{2}{5}$.}{}
%% TODO Darstellung für schwach fehlt in der Bank (brueche-dezimalzahlen-e2-k1-s7-v1; Abschnitt „Für schwache Schüler“)
\swz{Schreibe $\frac{13}{5}$ als gemischte Zahl.}{}
\swz[3]{Welcher Anteil der Figur ist grau? Gib ihn als vollständig gekürzten Bruch und in Prozent an. \hfill (P10 2014 OS)}{\bruchrechteck{12}{16}}
\end{aufgabe}

%% TODO Titel: Ich-kann-Satz fehlt in der Bank (Platzhalter: Kettenname) – Nr. 15
%% TODO Anweisung: ein Satz über der ersten Teilaufgabe fehlt in der Bank – Nr. 15
\begin{aufgabe}{gleichwertige Brüche am Streifen finden}
\swa{Der obere Streifen zeigt $\frac{1}{2}$. Wie viele Teile musst du im unteren Streifen färben, damit gleich viel gefärbt ist?}{\bruchrechteck{1}{2} \\ \bruchrechteck{0}{8}}
\end{aufgabe}

%% TODO Titel: Ich-kann-Satz fehlt in der Bank (Platzhalter: Kettenname) – Nr. 16
%% TODO Anweisung: ein Satz über der ersten Teilaufgabe fehlt in der Bank – Nr. 16
\begin{aufgabe}{Bruch als Geteilt-Aufgabe}
%% TODO Darstellung für schwach fehlt in der Bank (brueche-dezimalzahlen-e2-k3-s1-v1; Abschnitt „Für schwache Schüler“)
\swz{Schreibe $7 : 9$ als Bruch.}{}
\end{aufgabe}

%% TODO Titel: Ich-kann-Satz fehlt in der Bank (Platzhalter: Kettenname) – Nr. 17
%% TODO Anweisung: ein Satz über der ersten Teilaufgabe fehlt in der Bank – Nr. 17
\begin{aufgabe}{Kürzen und Erweitern – Fehler finden, Begründen, Darstellungswechsel, Anwendung}
\swz[3]{Mia erweitert $\frac{2}{7}$ mit $3$ und schreibt $\frac{6}{7}$. Finde den Fehler.}{}
\swfrage{Erweitert man $\frac{1}{2}$ mit $3$, stehen größere Zahlen da. Ist der Bruch größer geworden? Begründe.}
\swz{Jedes Teil des Streifens wird in $3$ gleiche Teile geteilt. Welcher Bruch steht dann für den gefärbten Anteil?}{\bruchrechteck{2}{5}}
\swz[3]{Eine Busfahrt dauert $45$ Minuten. Welcher Bruchteil einer Stunde ist das? Gib ihn gekürzt an.}{}
\end{aufgabe}

\uebersichtskasten{Erweitern: Zähler und Nenner mit derselben Zahl malnehmen – der Bruch bleibt gleich groß, die Teile werden feiner. \quad 2/3 = 4/6 = 8/12 \\ Kürzen: Zähler und Nenner durch dieselbe Zahl teilen. \quad 12/18 = 2/3 (durch 6) \\ Gleichnamig machen: beide Brüche auf denselben Nenner erweitern. \quad 1/4 und 5/6 → 3/12 und 10/12 \\ Gemischte Zahl: Ganze abspalten. \quad 11/4 = 2 3/4 \quad Bruch als Geteilt-Aufgabe: 5/8 = 5 : 8}
